\documentclass[12pt,a4paper]{article}
\usepackage[utf8]{inputenc}
\usepackage[russian,english]{babel}
\usepackage{amsmath, amssymb, amsthm}
\usepackage{graphicx}
\usepackage{hyperref}
\usepackage{cleveref}

\newtheorem{theorem}{Theorem}
\newtheorem{lemma}{Lemma}
\newtheorem{definition}{Definition}
\newtheorem{axiom}{Axiom}

\title{General Resonance Architecture (GRA): \\ From 6200 Daily Thoughts to the Nullification Game}
\author{Oleg Bitsoev (qqewq) \\ \texttt{https://orcid.org/0009-0004-1872-1153}}
\date{June 2026}

\begin{document}

\maketitle

\begin{abstract}
We present a unified mathematical framework --- the General Resonance Architecture (GRA) --- that models cognitive, social and artificial systems as hierarchical structures of resonant modes and internal contradictions (``foam''). The theory is anchored in the empirical estimate of ~6200 thoughts per day (Queen's University, 2020), each thought representing a micro-decision that shapes subjective reality. GRA introduces a hierarchy of abstraction levels, a foam functional $\Phi(k)$, stability criteria via derivatives with respect to hierarchy, and a nullification operator that resets destabilising resonances. We formalise the hierarchical stability rank $R_{\text{hier}}$ and show its application to personal coaching, multi-agent swarms, governance simulation, and even artistic creation. Verification strategies, backend/frontend implementations, and cross-disciplinary applications are outlined.
\end{abstract}

\section{Introduction}

The human mind produces a continuous stream of thoughts. Recent work by Queen's University (2020) estimated this stream at approximately 6200 thoughts per day, each thought corresponding to a transition between mental states. These first reactions to stimuli constitute the raw material from which our subjective reality is constructed. If we can learn to detect, structure and selectively ``nullify'' patterns in this stream, we can reshape not only behaviour but the very perception of the world.

The General Resonance Architecture (GRA) provides a formal language for this process. Originally developed for reducing ``argumentative foam'' in multi-agent debates, GRA has evolved into a broad framework applicable to individuals, organisations, swarms of AI agents, and even physical systems. This paper unifies the key concepts from over 100 GRA repositories into a coherent manifesto.

\section{Scientific Anchor: ~6200 Thoughts}

We adopt the estimate of ~6200 thoughts per day as a baseline for the daily flow of micro-decisions. For a subject $S$, let $M(t)$ denote the mental state at time $t$, with transitions $\Delta M_i = M(t_{i+1}) - M(t_i)$ for $i=1,\dots,N$, $N \approx 6200$. Each transition can be annotated with a meta-event $m_i$: content, emotion, decision.

\section{Hierarchical Model and Foam Functional}

A system is organised into levels $k = 0, 1, 2, \dots$ representing abstraction/control depth:
\begin{itemize}
\item $k=0$: raw sensations, actions
\item $k=1$: emotions, local plans
\item $k=2$: strategic goals, values
\item $k \gg 2$: meta‑agents, philosophy, meaning
\end{itemize}

For each level, define a foam functional $\Phi(k)$ as the expected mismatch, stress, or contradiction at that level. A simple analytic form is:
\begin{equation}
\Phi(k) = A e^{-\alpha k} + B e^{\beta k}, \quad A,B,\alpha,\beta > 0,
\end{equation}
where the first term captures stabilisation through abstraction, and the second term captures the risk of over-complexity.

\section{Axioms of GRA Nullification}

\begin{axiom}[Hierarchy]
The system is organised into levels $l=0,\dots,L$, each with state $x_l \in \mathbb{R}^{n_l}$.
\end{axiom}

\begin{axiom}[Foam]
For each level, a foam functional $\Phi(x_l) \ge 0$ is defined.
\end{axiom}

\begin{axiom}[Nullification]
There exists an operation $T_l$ such that $\Phi(T_l(x_l)) \le \Phi(x_l)$.
\end{axiom}

\begin{axiom}[Monotonicity]
In the ideal regime, $\Phi_{l+1} \le \Phi_l$.
\end{axiom}

\begin{axiom}[Sensitivity]
$\Phi(h)$ is differentiable w.r.t. the hierarchy parameter $h$; the sign of $d\Phi/dh$ determines evolution mode (productive collapse, stagnation, degradation).
\end{axiom}

\section{Stability Criteria}

\subsection{Continuous Criterion}
A nullified state at hierarchy level $h^*$ is \emph{stable} if:
\begin{equation}
\Phi(h^*) = 0,\qquad \frac{d\Phi}{dh}(h^*) = 0,\qquad \frac{d^2\Phi}{dh^2}(h^*) > 0.
\end{equation}
This ensures a local minimum of foam.

\subsection{Discrete Criterion}
For discrete levels $l$, the stable condition is:
\begin{equation}
\Phi_{l^*} = 0,\quad \Phi_{l^*+1} - \Phi_{l^*} = 0,\quad \Phi_{l^*+1} - 2\Phi_{l^*} + \Phi_{l^*-1} > 0.
\end{equation}

\subsection{Hierarchical Stability Rank $R_{\text{hier}}$}
Define the $n$-th finite difference magnitude:
\begin{equation}
D_n(k) = |\Delta^n \Phi(k)|,\qquad \Delta \Phi(k) = \Phi(k+1) - \Phi(k).
\end{equation}
The rank $R_{\text{hier}} = N$ is the minimal $n$ such that:
\begin{itemize}
\item For all $n < N$, $D_n(k)$ is monotonically decreasing for large $k$ (foam damping).
\item For $n = N$, $D_N(k)$ eventually stops decreasing and starts to grow or change sign (over‑hierarchy).
\end{itemize}
This rank characterises the system's optimal depth of reflection.

\section{Wave--Particle Duality in GRA}

Following the multiverse approach, we define a wavefunction over hierarchy levels:
\begin{equation}
\psi(k) \propto \exp(-\lambda \Phi(k)),\quad \lambda > 0.
\end{equation}
Then $|\psi(k)|^2$ is the probability that the system ``lives'' at level $k$. Sharp minima correspond to particle‑like stable goals; flat or multi‑modal profiles correspond to wave‑like wandering.

For two conflicting attractors (slits) $x_1^*$ and $x_2^*$, with foams $\Phi_1(k)$ and $\Phi_2(k)$, the total amplitude is:
\begin{equation}
A(k) = A_1 \psi_1(k) + A_2 \psi_2(k) e^{i \varphi(k)},
\end{equation}
where $\varphi(k) \sim \Phi_1(k) - \Phi_2(k)$. The interference term:
\begin{equation}
P(k) \propto |A(k)|^2 = |A_1|^2|\psi_1|^2 + |A_2|^2|\psi_2|^2 + 2\,\Re\{ A_1 A_2^* \psi_1 \psi_2^* e^{i\varphi} \}.
\end{equation}
A hard observation (projection) collapses the wave to a classical particle choice.

\section{Verification Frameworks}

\subsection{Mathematical Proofs}
The stability theorem is proven by showing that the conditions define a strict local minimum of $\Phi$. The swarm subjectivity theorem (GRA-Swarm-Subject) can be formalised using graph connectivity and convexity of the total foam functional.

\subsection{Simulation Benchmarks}
A benchmark suite (e.g., `GRA-Validation-Suite`) generates synthetic data with known ground truth and measures:
- Convergence speed to $\Phi \to 0$
- Recovery time after perturbation
- Scalability with number of agents/hierarchy levels

\subsection{Empirical Human Studies}
Design a mobile diary app (Angel‑in‑Pocket) to collect thought logs and emotional ratings. Compute $\Phi(k)$ via semantic embeddings and estimate $R_{\text{hier}}$ per user. Correlate with standard psychological scales (e.g., BDI, STAI) to validate predictive power.

\section{Implementation Blueprint}

\subsection{Backend (Python + FastAPI)}
- Endpoints: `/capture` (thought), `/foam` (profile), `/nullify` (reset), `/report`
- PostgreSQL for persistent storage, Redis for caching
- GRA core modules packaged as Python wheels

\subsection{Frontend (React + React Native)}
- Dashboard visualising $\Phi(k)$ as a hierarchical tree
- Real‑time thought capture via chat interface
- Interactive ``what‑if'' simulations (change weights, add new levels)

\section{Applications Across Disciplines}

\subsection{Natural Sciences}
- \textbf{Physics}: plasma control, quantum analogs (chiral nullification)
- \textbf{Biology}: immune system modelling (living vaccines), hormonal cycles
- \textbf{Chemistry}: materials discovery via multiverse hypothesis testing

\subsection{Humanities}
- \textbf{Psychology}: cognitive bias detection and therapy
- \textbf{Sociology}: policy simulation, conflict resolution
- \textbf{Linguistics}: semantic clustering, dialogue coherence

\subsection{Arts}
- \textbf{Music}: harmonic progression generation based on resonant modes
- \textbf{Painting}: colour and composition optimisation
- \textbf{Literature}: plot development, poetry generation (Ramanujan priors)

\section{Conclusion}

The GRA framework transforms the chaotic stream of ~6200 daily thoughts into a structured hierarchy of resonances and nullifications. By computing the stability rank $R_{\text{hier}}$ and applying the nullification operator, individuals, organisations and AI systems can achieve coherent, adaptive states. The theory is mathematically grounded, verifiable, and implementable, with a wide range of practical applications.

\section*{Acknowledgements}

The author thanks the open‑source community and all contributors to the GRA ecosystem.

\bibliographystyle{plain}
\begin{thebibliography}{9}

\bibitem{queens2020}
Queen's University (2020). ``Thought worms: segmenting brain activity into thought transitions.'' Unpublished preprint.

\bibitem{gra-hierarchical}
O. Bitsoev. \textit{GRA-Hierarchical-Stability: A Variational Formulation}. GitHub repository, 2026.

\bibitem{gra-rank}
O. Bitsoev. \textit{Hierarchical-Stability-Rank-N}. GitHub repository, 2026.

\bibitem{gra-swarm}
O. Bitsoev. \textit{GRA-Swarm-Subject}. GitHub repository, 2026.

\bibitem{gra-love}
O. Bitsoev. \textit{GRA-Love-Oriented-Nullification}. GitHub repository, 2026.

\bibitem{gra-unified}
O. Bitsoev. \textit{GRA-Unified}. GitHub repository, 2026.

\end{thebibliography}

\end{document}
